% SHARD-ID: RC-08C-WEIL-POSITIVITY-CONTINUOUS
% SHARD-TITLE: Weil Positivity, Continued: Continuous Time, Hilbert--P\'olya, the Zeta Entry, Prior Art
% SHARD-SUMMARY: Continuous (cMPS/Lindblad) version of Weil positivity via Bochner with the Dyson-expansion ring norms and the exact line-duality pairing identity.
% SHARD-SUMMARY: The Hilbert-Polya inner product exists iff the retained part is semisimple on the circle or line; Weil positivity is blind to Jordan blocks; the zeta instance as a conditional dictionary; prior art (Weil, Li, Huang, Suzuki, Hastings, Herglotz, Bochner, PT symmetry); the orchestrator's two observations and the numerics.
% SHARD-KEYWORDS: Bochner, Lorentzian, Dyson expansion, cMPS, Lindblad, Hilbert-Polya, Jordan block, Weil criterion, Li criterion, Huang criterion, screw function, PT symmetry, Kraus dichotomy
% SHARD-NOTES: none
% SHARD-SCRIPTS: scripts/weil_positivity.py

\section{Weil Positivity, Continued: Continuous Time, Hilbert--P\'olya, the Zeta Entry, Prior Art}
\label{sec:weil-positivity-continuous}

This section continues \cref{sec:weil-positivity} with the same objects
(\cref{def:continuous-rescaled-trace}); the status paragraph there applies.

\subsection{Continuous time: cMPS and Lindblad generators}

\begin{theorem}[Continuous Weil positivity is a half-plane bound]
\label{thm:weil-positivity-continuous}
\claimstatus{thm:weil-positivity-continuous}{proved}
For a linear operator $\Lgen$ on a finite-dimensional space, a trivial set $\trivS$ and
$\alev\in\mathbb R$, the rescaled continuous trace $\Fresc$ of
\cref{def:continuous-rescaled-trace} is continuous on $\mathbb R$ and positive definite iff
$\re\lmode\le\alev$ for every retained $\lmode$. A retained mode with $\re(\lmode-\alev) = -b<0$
and $\im\lmode = \omega$ contributes the Lorentzian $2b/(b^2+(x-\omega)^2)$ to the
spectral measure, a mode on the line a point mass, and positive definiteness forces
$|\Fresc|\le\Fresc(0)$, which a mode to the right of the line violates by continuous
mean-square growth.
\end{theorem}

\begin{theorem}[Line duality and the exact continuous pairing]
\label{thm:weil-line-duality}
\claimstatus{thm:weil-line-duality}{proved}
If the retained multiset is invariant under $\Rrefl{\alev}(\lmode) = 2\alev-\bar\lmode$, then
for $g = \sum_jc_j\delta_{\tsg_j}$ and $\hat g_{\Lgen}(\lmode) = \sum_jc_je^{\tsg_j(\lmode-\alev)}$,
\[
  \sum_{j,k}c_j\bar c_k\,\Fresc(\tsg_j-\tsg_k)
  \;=\; \sum_{\lmode\in\retA}\hat g_{\Lgen}(\lmode)\,\overline{\hat g_{\Lgen}(2\alev-\bar\lmode)} ,
\]
and $\Fresc$ is positive definite iff $\re\lmode = \alev$ for every retained $\lmode$.
\end{theorem}

This is the continuous inflow identity: the autocorrelation of the rescaled trace equals
the pairing of each mode with its reflection in the line, Weil's form in the form
$\sum_\lmode\hat g(\lmode)\overline{\hat g(1-\bar\lmode)}$.

\begin{theorem}[Dyson expansion as continuous ring norms]
\label{thm:cmps-ring-norms}
\claimstatus{thm:cmps-ring-norms}{proved}
For $\Lgen(x) = Kx+xK^\dagger+\sum_jR_jxR_j^\dagger$ on $\Mn$, with no trace
preservation assumed,
\[
  \Tr\nolimits_{\Mn}e^{\tsg\Lgen} \;=\; \sum_{k\ge0}\sum_{j_1,\dots,j_k}
  \int_{0<\tsg_1<\dots<\tsg_k<\tsg}
  \Big|\Tr\big(e^{(\tsg-\tsg_k)K}R_{j_k}e^{(\tsg_k-\tsg_{k-1})K}\cdots R_{j_1}e^{\tsg_1K}\big)\Big|^2
  \,d\tsg_1\cdots d\tsg_k \;\ge\;0,
\]
absolutely convergent, the $k = 0$ term being $|\Tr e^{\tsg K}|^2$. Moreover $\Lgen$
commutes with $x\mapsto x^\dagger$, so $\operatorname{spec}(\Lgen)$ is conjugation-closed.
\end{theorem}

The free propagators between jumps were missing from the draft; without them the product
is wrong for noncommuting $K$ and $R_j$. This is the continuous side A: a cMPS with
matrices $K$, $R_j$ has ring norms given by this sum, and the transfer generator $\Lgen$
is its side B. For a Hamiltonian generator $\Lgen = -i[H,\cdot]$ with $H = H^\dagger$ every
eigenvalue is purely imaginary and the line statement with $\alev = 0$ is trivial; that is
the Hilbert--P\'olya situation, and the content of RH in this language is the existence of
a presentation in which the transfer generator is a scalar plus $i$ times a Hermitian
operator, which the next two propositions make precise.

\subsection{The Hilbert--P\'olya inner product, and what positivity does not see}

\begin{proposition}[Unitarisability on the retained space]
\label{prop:hp-inner-product-discrete}
\claimstatus{prop:hp-inner-product-discrete}{proved}
Let $\trivS$ be a union of full eigenvalue classes, $V'$ the sum of the generalised
eigenspaces of the retained classes, $X' = X|_{V'}$. Then $X'$ is diagonalisable with all
eigenvalues of modulus $\critr$ iff there is an inner product on $V'$ in which $X'/\critr$
is unitary.
\end{proposition}

\begin{proposition}[Skew-adjoint realisation]
\label{prop:hp-inner-product-continuous}
\claimstatus{prop:hp-inner-product-continuous}{proved}
With the same convention, $\Lgen'-\alev$ is skew-adjoint in some inner product on $V'$ iff
$\Lgen'$ is diagonalisable with all eigenvalues on $\re\lmode = \alev$; equivalently $\Lgen'
= \alev+iH$ with $H$ Hermitian in that inner product.
\end{proposition}

\begin{proposition}[Weil forms do not see Jordan blocks]
\label{prop:weil-blind-jordan}
\claimstatus{prop:weil-blind-jordan}{proved}
For the Jordan block $X = \critr\left(\begin{smallmatrix}1&1\\0&1\end{smallmatrix}\right)$
with $\trivS = \varnothing$, $\nuseq{\wl} = 2$ for all $\wl$ and the sequence is positive
definite, although $X/\critr$ is not unitarisable; likewise $\Lgen = \alev+
\left(\begin{smallmatrix}0&1\\0&0\end{smallmatrix}\right)$. Weil positivity holds with
multiplicities and is blind to Jordan structure; a Hilbert--P\'olya inner product needs
semisimplicity in addition.
\end{proposition}

This separates two statements that are usually run together. ``All modes on the critical
circle'' is a statement about the characteristic polynomial and is equivalent to Weil
positivity. ``There is a Hilbert space in which the transfer operator is $\critr$ times a
unitary'' is stronger by exactly the absence of Jordan blocks. For the Riemann channel the
distinction is the multiplicity of the zeros: a multiple zero of $\Ssym$ gives a Jordan
block of $\Zt$ on $\KS$.

\subsection{The zeta instance and the prior art}

\begin{observation}[The zeta dictionary entry]
\label{obs:weil-zeta-dictionary}
\claimstatus{obs:weil-zeta-dictionary}{sketched}
For the compressed semigroup $\Zt$ with modes $\eta = -\bar\zero/2$, the critical real
part is $\alev = -\frac14$, the reflection $\zero\mapsto1-\bar\zero$ becomes
$\eta\mapsto-\frac12-\bar\eta$, the transform is $\hat g_\zeta(s) = \int g(\tsg)
e^{(s-1/2)\tsg/2}d\tsg$, and Weil's form $\sum_\zero\hat g(\zero)\overline{\hat g(1-\bar\zero)}$
is the autocorrelation of the centred trace distribution, whose full-line extension is
$\mathrm{T}_Z(-u) = e^{u/2}\overline{\mathrm{T}_Z(u)}$. The trivial modes at $0$ and $-\frac12$ and the
Gamma-factor terms are not eigenmodes of $\Zt$; they stay on the geometric side of the
explicit formula. The passage from finitely many modes to $\zeta$ needs the completed-zeta
data, a distributional trace realisation, the full explicit formula and the infinite Weil
converse, with Bochner--Schwartz in place of Bochner.
\end{observation}

The prover's T6 also derives that in the explicit-formula convention of
\cref{sec:riemann-channel} the prime part of the centred trace carries weight
$-2\vM(n)n^{-1/2}$ at $\pm2\log n$; this is \cref{obs:ringnorm-sign-correction}.

\begin{citedfact}[Weil's criterion]
\label{cit:weil-criterion}
\claimstatus{cit:weil-criterion}{cited}
RH holds iff the Weil scalar product $\sum_\zero F(\zero)\overline{G(1-\bar\zero)}$ is
positive on the admissible test class \cite{connesconsani2020weil,lagarias2005li}.

\par\noindent Quoted (\texttt{prov:cc20-weilcrit}, \texttt{prov:lagarias05-weil-scalar-product}):
{\small\texttt{\detokenize{In fact, following \cite{yoshida}, it is enough to prove the negativity of the right-hand side of \eqref{explicit form} for functions $g$ with compact support in the (locally compact) multiplicative group $\R_+^*=(0,\infty)$. Furthermore, given any finite set of complex numbers $F\supset \{0,1\}$, $F\cap Z=\emptyset$, using the notation $g^\sharp(x):=x^{-1}g(x^{-1}), \ \bar g(x)=\overline{g(x)}$, one has (Appendix \ref{apppositivity}, Proposition~\ref{mainprop}) \begin{equation}\label{weilnegav} RH \iff \sum_v {\mathcal W}_v(g*\bar g^\sharp)\leq 0, \quad \forall g\in C_c^\infty(\R_+^*)\mid \tilde g(z)=0,\ \forall z\in F. \end{equation}}}}

{\small\texttt{\detokenize{Given $F(s), G(s) \in \sA$ we define the {\em Weil scalar product} associated to the automorphic representation $\pi$ by \beql{301} \langle F, G \rangle_{\sW(\pi)} := \sum_{ \rho \in Z(\pi) } F(\rho) \overline{G( 1 - \bar{\rho})}. \eeq The sum on the right counts zeros with}}}
\end{citedfact}

\begin{citedfact}[Li's criterion]
\label{cit:li-criterion}
\claimstatus{cit:li-criterion}{cited}
RH iff $\lambda_n\ge0$ for all $n\ge1$; each condition $\lambda_n\ge0$ encodes Weil
positivity for one test function \cite{voros2006li,lagarias2005li}.

\par\noindent Quoted (\texttt{prov:voros06-li-criterion}, \texttt{prov:lagarias05-li-is-weil}):
{\small\texttt{\detokenize{\emph{Li's criterion} for the Riemann Hypothesis (RH) states that the latter is true if and only if a specific real sequence $\{ \lambda_n \}_{n=1,2,\ldots}$ has \emph{all its terms positive} \cite{LI1,BL}.}}}

{\small\texttt{\detokenize{\begin{theorem}~\label{th31} Let $\pi$ be an irreducible cuspidal unitary automorphic representation of $GL(N)$, with associated $\xi$-function $\xi(s, \pi)$ and with Weil scalar product $\langle \cdot, \cdot \rangle_{\sW(\pi)}$. For the Li test functions $G_n(s) = 1 - (1 - \frac{1}{s})^n$ there holds \beql{304} \langle G_n, G_m \rangle_{\sW(\pi)} = \lambda_n(\pi) + \lambda_{-m}(\pi) - \lambda_{n-m}(\pi). \eeq}}}
\end{citedfact}

\begin{citedfact}[Huang's criterion for regular graphs]
\label{cit:huang-graph-criterion}
\claimstatus{cit:huang-graph-criterion}{cited}
A connected $(\qs+1)$-regular graph on $n$ vertices is Ramanujan iff $h_k\ge0$ for all
$k\ge1$, where for odd $k$, $h_k = 2(n-1)+\qs^{k/2}+\qs^{-k/2}-\qs^{-k/2}N_k$ with $N_k$ the
number of cyclically non-backtracking closed walks of length $k$ \cite{huang2019ihara}.

\par\noindent Quoted (\texttt{prov:huang19-ramanujan-iff-positive}, \texttt{prov:huang19-hk-cycle-formula}):
{\small\texttt{\detokenize{\begin{thm}\label{thm:Li} The following are equivalent: \begin{enumerate} \item $X$ is Ramanujan. \item $h_k\geq 0$ for all $k\geq 1$. \item $h_k\geq 0$ for infinitely many even $k\geq 2$. \end{enumerate} \end{thm}}}}

{\small\texttt{\detokenize{2(n-1) + q^{\frac{k}{2}} +q^{-\frac{k}{2}} -q^{-\frac{k}{2}} N_k \qquad &\hbox{for all odd $k$},}}}
\end{citedfact}

\begin{citedfact}[Hastings' bound]
\label{cit:hastings-bound}
\claimstatus{cit:hastings-bound}{cited}
For a quantum expander with $\Dk$ unitary Kraus operators the Ramanujan value of the
second largest eigenvalue modulus is $\lambda_H = 2\sqrt{\Dk-1}/\Dk$ \cite{hastings2007random}.

\par\noindent Quoted (\texttt{prov:hastings07-lambdaH}):
{\small\texttt{\detokenize{Let $\lambda_2$ be the eigenvalue with the second largest absolute value of all eigenvalues other than $\lambda_1$. Let \be \lambda_{H}=\frac{2\sqrt{D-1}}{D}. \ee The main result of this paper is that, in the Hermitian case, for any $\epsilon>0$ the probability that $|\lambda_2|$ is within $\epsilon$ of $\lambda_{H}$ approaches unity as $N\rightarrow \infty$. Interestingly, this is the same as the recently proven tight bound\cite{fried} in the classical case, but the proof in the quantum case is much simpler.}}}
\end{citedfact}

\begin{citedfact}[Herglotz and Bochner]
\label{cit:herglotz-bochner}
\claimstatus{cit:herglotz-bochner}{cited}
A sequence is positive definite iff its Toeplitz matrices are positive semidefinite iff it
is the Fourier coefficient sequence of a positive measure on the circle
\cite{alpay2014herglotz}; a continuous positive definite function on $\mathbb R$ is the
Fourier transform of a finite positive measure \cite{norvidas2020bochner}.

\par\noindent Quoted (\texttt{prov:acks14-herglotz}, \texttt{prov:norvidas20-bochner}):
{\small\texttt{\detokenize{A result of Herglotz, also known as Bochner's theorem, asserts that: \begin{Tm} \label{Tm:Oct27yt1} The function $n \mapsto r(n)$ from $\mathbb Z$ into $\mathbb C^{s \times s}$ is positive definite if and only if there exists a unique positive $\mathbb C^{s \times s}$-valued measure $\mu$ on $[0, 2\pi ]$ such that \begin{equation} \label{eq:Oct27nv1} r(n) = \int_0^{2\pi} e^{i n t} d\mu(t), \quad n \in \mathbb Z. \end{equation} \end{Tm}}}}

{\small\texttt{\detokenize{A function $f:\R\to\Co$ is said to be positive definite if \begin{equation}\label{1.1} \sum_{j, k=1}^n f(x_j-x_k)c_j{\overline{c}}_k\ge 0 \end{equation} holds for all finite sets of complex numbers $c_1,\dots, c_n$ and points $x_1,\dots,x_n\in\R$. A survey about positive definite functions and its generalizations can be found in \cite{9}. The Bochner theorem states that a continuous function $f$ on $\R$ is positive definite if and only if it is the Fourier transform of a finite nonnegative measure $\mu$ on $\R$, i.e. \[ f(x)= \hat{\mu}(x)=\int_{-\infty}^{\infty}e^{-ixt}\,d\mu(t), \]}}}
\end{citedfact}

\begin{citedfact}[Suzuki's kernel form]
\label{cit:suzuki-kernel}
\claimstatus{cit:suzuki-kernel}{cited}
RH is equivalent to nonnegative definiteness of the kernel $G(t,u) =
g(t-u)-g(t)-g(-u)+g(0)$ built from a screw function of $\zeta$, in the
Herglotz--Nevanlinna framework, with Li's criterion as a case of Weil positivity
\cite{suzuki2022screw}.

\par\noindent Quoted (\texttt{prov:suzuki22-rh-iff-nnd}, \texttt{prov:suzuki22-weil-positivity}):
{\small\texttt{\detokenize{\begin{theorem} \label{Thm_1_3} The RH is true if and only if the hermitian form $\langle \cdot,\cdot \rangle_{G_g,a}$ is non-negative definite on $\mathfrak{C}_0(a)$, that is, $\langle \phi,\phi \rangle_{G_g,a} \geq 0$ for all $\phi \in \mathfrak{C}_0(a)$ for every $0<a<\infty$. Moreover, assuming that the RH is true, $\langle \cdot,\cdot \rangle_{G_g,a}$ is positive definite on $L^2(-a,a)$, that is, $\langle \phi,\phi \rangle_{G_g,a}>0$ for all non-zero $\phi \in L^2(-a,a)$ for every $0<a<\infty$. \end{theorem}}}}

{\small\texttt{\detokenize{showed that the RH is true if and only if the distribution $W$ is non-negative definite, that is,}}}
\end{citedfact}

\begin{citedfact}[PT symmetry and real spectrum]
\label{cit:pt-real-spectrum}
\claimstatus{cit:pt-real-spectrum}{cited}
A PT-symmetric non-Hermitian Hamiltonian has real spectrum when the symmetry is unbroken,
i.e.\ its eigenstates are also PT eigenstates \cite{bender2017hamiltonian}.

\par\noindent Quoted (\texttt{prov:bbm17-pt-real}):
{\small\texttt{\detokenize{(iii) Although $\hat H$ is not Hermitian, $\ri{\hat H}$ is $\cPT$ symmetric; that is, $\ri{\hat H}$ is invariant under parity-time reflection (in the sense to be defined), which means that the eigenvalues of $\ri{\hat H}$ are either real or else occur in complex-conjugate pairs. If $\ri{\hat H}$ has {\it maximally broken} $\cPT$ symmetry; that is, if all of its eigenvalues are {\it pure-imaginary complex-conjugate pairs}, then the eigenvalues of ${\hat H}$ are real and the Riemann hypothesis follows. (iv) While ${\hat H}$ is not Hermitian (symmetric)}}}
\end{citedfact}

Positivity criteria for RH are old (Weil; Li; the kernel form of Suzuki), and Huang's
criterion is the graph case. What is not in the literature found is the statement for an
arbitrary transfer operator, the separation of the one-sided bound from the duality, and
the Kraus dichotomy. Huang's criterion is termwise rather than Toeplitz; the next
observation reconciles the two.

\subsection{Two observations by the orchestrator}

\begin{observation}[Huang's criterion is the boundedness form]
\label{obs:huang-boundedness}
\claimstatus{obs:huang-boundedness}{proved}
In finite dimensions, for a conjugation-closed retained multiset, the following are
equivalent: positive definiteness of $\nuseq{\wl}$; boundedness; $|\nuseq{\wl}|\le\nuseq{0}$
for all $\wl$; and the termwise criterion $\re\nuseq{\wl}\le\nuseq{0}$ for all $\wl\ge1$.
For a connected non-bipartite $(\qs+1)$-regular graph with $\critr = \sqrt{\qs}$ and
$\trivS = \{+1^{[m-n]},-1^{[m-n]},\qs,1\}$ ($m$ edges, $n$ vertices), Huang's sequence is
exactly $h_k = \nuseq{0}-\nuseq{k}$ with $\nuseq{0} = 2(n-1)$, so Huang's criterion is the
boundedness form of Weil positivity. The Toeplitz form is what survives when the trace is
a distribution and boundedness has no meaning.
\end{observation}

\begin{observation}[The Kraus dichotomy]
\label{obs:kraus-dichotomy}
\claimstatus{obs:kraus-dichotomy}{proved}
For any $\Ad$-type family the rings are positive and $\operatorname{spec}(\Sig)$ and
$\operatorname{spec}(\Hash)$ are conjugation-closed. (b) Inverse pairing buys the functional
equation: with $\trivS = \{+1^{[k]},-1^{[k]}\}$ and $\critr = \sqrt{\Dk-1}$, ``all retained
modes on the circle'' is equivalent to: every eigenvalue $\sigev$ of $\Sig$ is real with
$|\sigev|\le2\sqrt{\Dk-1}$. (c) Adjoint pairing buys Hilbert--Schmidt self-adjointness of
$\Sig$, so $\operatorname{spec}(\Sig)$ is real for free; $\operatorname{spec}(\Hash)$ need
not be real, and in general there is no functional equation
(\cref{prop:kraus-no-duality-example}). (d) A family with both pairings is unitary; there
the eigenvalue $\sigev = \Dk$ of $\Sig$ (from $\chan(1) = 1$) always violates the band, so
(b) is applied with the enlarged trivial set of \cref{thm:kraus-weil-criterion}(ii),
$\trivS = \{+1^{[k]},-1^{[k]}\}\uplus\{\qs^{[d_+]},1^{[d_+]}\}\uplus\{(-\qs)^{[d_-]},(-1)^{[d_-]}\}$,
and reduces to Hastings' bound alone. Reality of $\operatorname{spec}(\Sig)$ is the
Hilbert--P\'olya half; the functional equation is the other half; a Kraus family gets the
two from two different pairings, and gets both only by being unitary.
\end{observation}

The dichotomy reopens \cref{conj:kraus-ramanujan} in a sharper form: inverse-paired
non-unitary families (for instance $B_i = GU_iG^{-1}$ with a common similarity $G$, which
keeps $\Sig$ similar to a self-adjoint operator) have a genuine two-sided Ramanujan
property with a positive side A, and their non-Hermitian $\Sig$ with real spectrum is the
transfer-matrix version of an unbroken PT symmetry (\cref{cit:pt-real-spectrum}).

\begin{numerical}[Numerical checks]
\label{num:weil-positivity}
\claimstatus{num:weil-positivity}{numerical}
\texttt{scripts/weil\_positivity.py} (\texttt{outputs/weil\_positivity.txt}): random
adjoint-paired families ($n = 2,3$; $\Dk = 4,6$; 300 cases) satisfy ``one-sided bound
$\Rightarrow$ Toeplitz positivity at $L = 40$'' without exception, the threshold radius
approaches $\max|\edgeev|$ off $\trivS$ from below like $L^{-3/2}$ (relative gap
$3\cdot10^{-2}$, $5\cdot10^{-3}$, $4.7\cdot10^{-4}$ at $L = 40, 160, 640$), the retained
spectra are conjugation-closed to $10^{-14}$ and not reciprocal-invariant; Haar-unitary
families reproduce \cref{cor:qihara-unitary} to $10^{-15}$, the $\reflJ$-invariance, and
``Weil positivity iff Hastings' bound'' on four passing and two violating families, with
$\Wform = \Mform$ to $10^{-14}$; an inverse-paired non-unitary family has reciprocal-invariant,
conjugation-closed spectrum and nonnegative ring traces to $10^{-15}$; random
Lindblad-type generators have $\Tr e^{\tsg\Lgen}\ge0$ matching the Dyson expansion to
$10^{-7}$ and Gram positivity iff $\max\re\lmode\le\alev$; the Jordan block passes with
$\nuseq{\wl}$ identical to its diagonal companion. Three of thirty checks fail: they are the
three drafted statements the prover corrected.
\end{numerical}

Two remarks on the numerics. Positivity at a fixed finite $L$ is weaker than the theorem:
a mode slightly outside the circle is masked by the positive contribution of the interior
modes until $L$ is large. And the $\Dk = 4$ Pauli family $\{X,Z,X^\dagger,Z^\dagger\}$ is an
exact Ramanujan example with $\chan$-spectrum $\{1,0,0,-1\}$, exercising the $-1$ branch
of the trivial set with retained modes $\pm i\sqrt3$.
